% ==============================================================================
% US-072 -- FarSLIP contrastivo-fenologico (seccion Method/Experiments modular).
%
% Seccion auto-contenida del Paper Track (E11). Documenta el diagnostico del bug
% torch.randn, el fix con prototipos fenologicos reales, el filtro 3:1 Meadow
% per-patch, la ablacion de bandas, el protocolo incremental 4->18 y el marco
% Be My Eyes. TODAS las cifras provienen de artefactos REALES del repo (ver los
% comentarios de cada \input{...}). Las tablas viven en
% paper/tables/farslip-method/ y se incluyen al final de cada subseccion.
%
% Para compilar de forma autonoma (fuera del manuscrito principal) se requiere:
%   \graphicspath{{../figures/}{../../reports/farslip/figures/}}
% y los \bibitem li2025farslip / wen2025phenology / huang2025bemyeyes
% (ver paper/bib/farslip_refs.bib). Si se incrusta en
% docs/final_doc/Avance7_equipo17.tex, esos recursos ya existen alli.
% ==============================================================================

\section{FarSLIP contrastivo-fenologico: del fallo contra ruido a la senal
visión-lenguaje}
\label{sec:farslip_method}

Esta seccion documenta el camino metodologico que eleva FarSLIP de una
ablacion negativa --- perdia F1-macro $0.163$ frente al $0.233$ de AlphaEarth
sobre las 18 clases de PASTIS-R --- a un componente visión-lenguaje cuyo aporte
es \textit{complementario y honesto}. El relato es deliberadamente transparente:
reportamos el bug que originaba el resultado negativo, el fix que cierra el gap
\textit{propio} de FarSLIP, y el hecho de que, aun corregido, el espacio FarSLIP
fiel \textbf{no supera} a AlphaEarth como clasificador denso. Su valor reside en
el ensamble heterogeneo y en la transferibilidad multi-region (Sec.~\ref{sec:farslip_transfer}).

% ------------------------------------------------------------------------------
\subsection{F1 --- Diagnostico del bug: alineamiento contrastivo contra ruido}
\label{sec:farslip_f1}

La perdida de alineamiento region-categoria de FarSLIP~\cite{li2025farslip}
(\texttt{RegionCategoryAlignmentLoss}, un InfoNCE con temperatura $0.07$) alinea
el token \texttt{[CLS]} del estudiante contra una matriz de \textit{prototipos
de texto}. En la implementacion original
(\texttt{ml/farslip/train.py}, ruta \texttt{proto\_source="random"}) esa matriz
se inicializaba con \texttt{torch.randn(n\_protos, hidden\_dim)}: ruido gaussiano
sin contenido semantico. El contraste empujaba el \texttt{[CLS]} hacia vectores
aleatorios, de modo que la senal categorica nunca se anclaba a una descripcion
real del cultivo. El docstring de \texttt{build\_text\_prototypes} nombra el
sintoma textualmente: \textit{``contrastive alignment against noise --- the bug
that made FarSLIP lose 0.163 vs 0.233''}. Este es el origen del resultado
negativo: no un fallo del modelo de fundacion, sino una inicializacion que
convertia una perdida contrastiva en regularizacion hacia ruido.

% ------------------------------------------------------------------------------
\subsection{F2 --- Prototipos fenologicos reales y reproyeccion ortogonal frozen}
\label{sec:farslip_f2}

El fix sustituye el ruido por \textit{prototipos fenologicos reales}, una clase
a la vez. Para cada una de las 18 clases de PASTIS-R se redacta una descripcion
fenologica siguiendo el marco de Wen et al.~\cite{wen2025phenology} (inicio de
verdor, pico de NDVI, fecha de senescencia, textura de dosel), generada con
Gemini~Flash y codificada con MiniLM (\texttt{all-MiniLM-L6-v2}) a un embedding
de 384 dimensiones L2-normalizado. El flujo (ruta \texttt{proto\_source="pastis"})
es:

\begin{enumerate}[leftmargin=*]
    \item \texttt{load\_class\_prototype\_embeddings()} $\rightarrow$
    \texttt{proto\_18} $(18, 384)$, prototipos MiniLM por clase.
    \item \texttt{expand\_to\_cap()} $\rightarrow$ \texttt{proto\_cap}
    $(32, 384)$, mapeo del esquema PASTIS-18 al vocabulario CAP de 32 categorias.
    \item \texttt{np.tile()} $\rightarrow$ \texttt{proto\_tiled}
    $(n\_regions \cdot 32, 384)$ en orden \textit{region-major}
    (\texttt{region $\cdot$ n\_categories + category}), coincidente con el target
    de la perdida.
\end{enumerate}

El \textit{tile} de 384 dimensiones se reproyecta a la dimension del
\texttt{[CLS]} del estudiante mediante una matriz semi-ortogonal $W$ generada
con \texttt{torch.nn.init.orthogonal\_}, \texttt{requires\_grad=False}: nunca
entra al optimizador ni al \texttt{state\_dict}. Al ser ortogonal preserva la
norma $L_2$ y los angulos del subespacio MiniLM, de modo que las 18 clases
permanecen distinguibles tras la reproyeccion.

\paragraph{Correccion factual sobre la dimensionalidad.}
La reproyeccion de los prototipos del \textit{loss} es \textbf{$384\rightarrow768$},
no $384\rightarrow512$: el \texttt{hidden\_size} del teacher CLIP ViT-B/16 es
$768$, confirmado por \texttt{trainer.teacher.config.hidden\_size} y por la
columna \texttt{n\_dims}$=768$ del artefacto de evaluacion. La cifra ``512-dim
por parcela'' corresponde a otro objeto --- el banco de \textit{features} que el
extractor produce a nivel de parcela --- y no debe confundirse con la
reproyeccion de prototipos del InfoNCE. Distinguimos explicitamente:
banco-parcela $512$-dim frente a prototipos/\texttt{[CLS]} $768$-dim.

\paragraph{Efecto del fix sobre el gap propio.}
El fix cierra el gap \textit{propio} de FarSLIP: el espacio fiel v2 pasa de
F1-macro $0.163$ (contra ruido) a $0.5551 \pm 0.0205$ (con prototipos
fenologicos) sobre las mismas 567 parcelas, una mejora de $+0.392$ respecto al
baseline buggeado. La Tabla~\ref{tab:farslip_vs_alphaearth} reporta la
comparacion honesta frente a AlphaEarth.

% Tabla US-072 -- FarSLIP fiel v2 vs AlphaEarth (separabilidad + F1 macro).
% Cifras EXACTAS de reports/farslip/metrics/us037_farslip_fiel_vs_alphaearth.csv
% (n_samples=567, 5-fold spatial CV). NO editar a mano: regenerar desde el CSV.
\begin{table}[htbp]
\caption{Espacio de representacion FarSLIP fiel (v2) frente a AlphaEarth sobre
las mismas 567 parcelas PASTIS-R. El espacio FarSLIP (768-dim) NO supera a
AlphaEarth (64-dim, Satellite Embedding V1 Annual, data version 1.1) ni en
separabilidad (silhouette) ni en F1-macro; se reporta el resultado honesto.}
\label{tab:farslip_vs_alphaearth}
\centering
\begin{tabular}{lccccc}
\toprule
\textbf{Espacio} & \textbf{Silhouette} & \textbf{F1-macro} & \textbf{$n$ dim} & \textbf{$n$ muestras} & \textbf{$\Delta$ vs AlphaEarth} \\
\midrule
FarSLIP-pheno (fiel v2) & 0.0120 & $0.5551 \pm 0.0205$ & 768 & 567 & $-0.0024$ \\
AlphaEarth (2019)       & 0.0145 & $0.6446 \pm 0.0415$ &  64 & 567 & ref. \\
\bottomrule
\end{tabular}
\end{table}


\paragraph{Claim defendible.}
Corregido el bug, el espacio FarSLIP fiel \textbf{sigue por debajo} de AlphaEarth
($-0.0024$ en F1-macro, $-0.0025$ en silhouette). No afirmamos que FarSLIP supere
a AlphaEarth ni a Gemini: el aporte de FarSLIP es como miembro complementario del
ensamble y en el transfer multi-region, no como clasificador denso ganador. Esta
es la lectura que el resto de la seccion sostiene.

% ------------------------------------------------------------------------------
\subsection{F3 --- Filtro 3:1 Meadow por parcela (\texttt{dominance\_ratio})}
\label{sec:farslip_f3}

PASTIS-R esta dominado por \textit{Meadow} (clase 1) y \textit{Background}
(clase 0), que juntos pueden cubrir $>80\%$ de un patch. Extraer
\textit{features} sobre esos patches contamina el espacio de embeddings con
senal poco informativa. El filtro previo de Isaac operaba en modo
\texttt{coverage} (legacy): conservaba un patch si la cobertura combinada de las
clases objetivo sobre los pixeles validos superaba un umbral global. El presente
trabajo introduce el modo \texttt{dominance\_ratio} en
\texttt{ml/data/pastis\_filter.py}: conserva un patch si
$\text{Meadow}_{px} \leq 3 \cdot \text{second}_{px}$, donde
$\text{second}_{px}$ es el conteo mayor \textit{en ese mismo patch} entre las
clases objetivo, excluyendo Meadow, Background~(0) y Void~(19) por identificador
literal (no via \texttt{ignore\_index}) para no contaminar la razon. La
comparacion es inclusiva ($\leq$): un patch exactamente 3:1 se conserva.

La diferencia con el filtro \texttt{coverage} es conceptual: \texttt{coverage}
mide pureza \textit{global} de clases objetivo; \texttt{dominance\_ratio} acota
el desbalance \textit{local} de Meadow respecto a la segunda clase de cada patch,
descartando solo los patches sobre-dominados por prado sin tirar la senal
minoritaria del resto. Casos borde explicitos y testeados:
$\text{Meadow}_{px}=0$ conserva (metrica $0.0$); $\text{Meadow}_{px}>0$ con
$\text{second}_{px}=0$ (sin clase competidora) descarta (metrica $\infty$).

% ------------------------------------------------------------------------------
\subsection{F4 --- Ablacion de bandas (rgb / nir\_rgb / 4band)}
\label{sec:farslip_f4}

Se evaluaron tres variantes de seleccion de bandas (US-035), implementadas como
\textit{transform} por recorte en \texttt{ml/farslip/bands.py}: \textit{rgb}
(B04-B03-B02, 3 canales), \textit{nir\_rgb} (B08-B04-B03 falso color, 3 canales)
y \textit{4band} (B02-B03-B04-B08 identidad, 4 canales). La variante 4band
reutiliza la adaptacion $3\rightarrow4$ del \texttt{patch\_embed} con
NIR\,=\,mean(RGB) (US-034) para evitar neuronas muertas al ampliar el primer
bloque convolucional.

% Tabla US-072 -- Ablacion de bandas FarSLIP (US-035). Cifras EXACTAS extraidas
% de los logs reales reports/farslip/logs/{baseline-rgb,baseline-nir,4band-pheno}.log
% (linea "training done", epoch 3, ultima de 4 epochs).
% CAVEAT HONESTO: los 3 logs solo registran las PERDIDAS de destilacion finales
% (loss_total, loss_patch, loss_cls, loss_aux); NO contienen la metrica de
% evaluacion consolidada (F1-macro/mIoU) por variante. El F1/mIoU por variante
% queda PENDIENTE (ver docs/blockers/epic11-notas.md, US-072). NO inventar la cifra.
\begin{table}[htbp]
\caption{Ablacion de bandas FarSLIP (US-035): tres corridas reales de
destilacion CLIP region-aware (4 epochs). \textit{rgb} = B04-B03-B02 (3 canales);
\textit{nir\_rgb} = B08-B04-B03 falso color (3 canales); \textit{4band} =
B02-B03-B04-B08 identidad (4 canales, adaptacion 3$\rightarrow$4 con
NIR\,=\,mean(RGB) anti-neurona-muerta, US-034). Se reportan las perdidas finales
de destilacion; la metrica de evaluacion (F1-macro/mIoU) por variante no quedo
consolidada en los logs y se marca como pendiente.}
\label{tab:band_ablation}
\centering
\begin{tabular}{lccccc}
\toprule
\textbf{Variante} & \textbf{Canales} & \textbf{loss\_total} & \textbf{loss\_patch} & \textbf{loss\_cls} & \textbf{F1-macro} \\
\midrule
rgb     & 3 & 2.2893 & 0.0179 & 4.5244 & pendiente \\
nir\_rgb & 3 & 2.2858 & 0.0181 & 4.5166 & pendiente \\
4band   & 4 & 2.3553 & 0.0786 & 4.5510 & pendiente \\
\bottomrule
\end{tabular}
\end{table}


\paragraph{Caveat de evaluacion (resultado parcial honesto).}
Las tres corridas existen y convergen en perdida de destilacion
(Tabla~\ref{tab:band_ablation}): la perdida total final es muy similar entre
\textit{rgb} ($2.2893$) y \textit{nir\_rgb} ($2.2858$), y algo mayor en
\textit{4band} ($2.3553$), con una perdida de \textit{patch} notablemente mas
alta en 4band ($0.0786$ frente a $\approx0.018$) coherente con el canal NIR
adicional. Sin embargo, los logs no registran la metrica de evaluacion
consolidada (F1-macro/mIoU) por variante; por tanto NO reportamos un ganador de
clasificacion entre bandas. Reportar una cifra de F1 inexistente seria inventar
datos. El cierre de esta tabla con la metrica de evaluacion queda como tarea
pendiente documentada en los \textit{blockers} de la epica.

% ------------------------------------------------------------------------------
\subsection{F5 --- Protocolo incremental 4$\rightarrow$18 y fusion en el ensamble}
\label{sec:farslip_f5}

El clasificador FarSLIP se construye de forma incremental
(\texttt{ml/farslip/incremental\_curriculum.py}): se parte de las 4 clases mas
cardinales (Stage-1) y se anaden 2 clases por paso hasta 18 (Stage-2), inicializando
cada etapa desde la anterior con \texttt{strict=False} para permitir el cambio de
dimension del cabezal. El criterio de parada (\texttt{stop\_criterion}) contempla
tres razones explicitas: \textit{(i)} clases nuevas inaceptables
(F1$<0.30$), \textit{(ii)} degradacion de clases previas
($\text{drop}>0.05$, olvido catastrofico), \textit{(iii)} maximo de clases
alcanzado.

La Tabla~\ref{tab:cardinality_sweep} muestra el barrido de cardinalidad por
parcela --- la evidencia empirica que motiva el protocolo. La dificultad crece
monotonicamente con el numero de clases: F1-macro $0.7025$ con 4 clases,
$0.4579$ con 6, $0.4075$ con 8, $0.3589$ con 10 y $0.3328$ con 12. El numero de
parcelas evaluables crece (de 1301 a 3200) al admitir clases mas raras, pero la
calidad promedio cae: las clases minoritarias con fenologia ambigua arrastran el
F1-macro.

% Tabla US-072 -- Barrido de cardinalidad por parcela (evidencia del protocolo
% incremental 4->18). Cifras EXACTAS de reports/farslip/metrics/parcel_sweep.csv.
% Cada fila es un checkpoint real entrenado SOLO sobre las K clases mas cardinales.
% NO editar a mano: regenerar desde el CSV.
\begin{table}[htbp]
\caption{Barrido de cardinalidad por parcela del clasificador FarSLIP-pheno.
La dificultad crece monotonicamente con el numero de clases: de 4 clases
($0.7025$) a 12 clases ($0.3328$). Evidencia empirica que motiva el protocolo
incremental 4$\rightarrow$18 (Stage-1$\rightarrow$Stage-2).}
\label{tab:cardinality_sweep}
\centering
\begin{tabular}{ccccc}
\toprule
\textbf{$K$ clases} & \textbf{F1-macro} & \textbf{mIoU} & \textbf{bien resueltas (F1$\geq$0.5)} & \textbf{$n$ parcelas eval} \\
\midrule
 4 & 0.7025 & 0.5547 & 4 & 1301 \\
 6 & 0.4579 & 0.3120 & 2 & 1843 \\
 8 & 0.4075 & 0.2657 & 3 & 2301 \\
10 & 0.3589 & 0.2288 & 3 & 2706 \\
12 & 0.3328 & 0.2071 & 1 & 3200 \\
\bottomrule
\end{tabular}
\end{table}


\paragraph{Fallback honesto a 18 clases desde cero.}
El baseline fiel de 18 clases sin pesos de clase (variante
\texttt{faithful\_v2}, \textit{region\_category} MPCL + $L_{glo}$) alcanza
F1-macro $0.164$ y mIoU $0.111$, con 4 clases de senal clara
(Meadow $0.761$, Orchard $0.516$, Grapevine $0.489$, Corn $0.464$) y 14 clases
en cero o muy bajo. La variante con pesos de clase \textit{inverse-frequency}
(\texttt{faithful\_v2\_cw}) EMPEORA el balance neto a F1-macro $0.102$: desinfla
los dominantes mas de lo que rescata a los raros. Estos resultados se reportan
sin maquillaje; las 18 clases son $\approx4.5\times$ mas dificiles que las 4
faciles ($0.576$ en US-036-a, no comparable directamente).

\paragraph{Fusion en el ensamble heterogeneo.}
El aporte neto de FarSLIP se mide en el meta-modelo de \textit{stacking}. El
ensamble de 5 miembros (TSViT-pheno, U-TAE, XGB-AlphaEarth, FarSLIP-ft18,
FarSLIP-zeroshot) sobre probabilidades \textit{out-of-fold} alcanza F1-macro
$0.6477$ (\textit{accuracy} $0.7935$) frente a $0.6359$ del ensamble de 3
miembros sin FarSLIP: un delta de $+0.0118$ atribuible a las dos ramas FarSLIP en
el \textit{stacking} OOF. El manuscrito principal reporta la cifra del ensamble
final consolidado (F1-macro $0.749$); la magnitud exacta de la mejora atribuible
a FarSLIP en esa cifra final difiere de este delta OOF y se trata como decision
editorial (ver \textit{blockers}). En todo caso, la conclusion es estable: el
aporte de FarSLIP es \textbf{positivo pero pequeno} ($\approx+0.3\%$ a $+1.2\%$
segun la metrica), nunca el motor del clasificador.

% ------------------------------------------------------------------------------
\subsection{Marco \textit{Be My Eyes}: el LLM no clasifica pixeles}
\label{sec:farslip_bemyeyes}

El resultado de que los prototipos fenologicos de texto aporten solo
$\approx0.3\%$ como clasificador supervisado --- y no el $5\%$ que una lectura
optimista sugeriria --- es un resultado \textit{esperado y correctamente
enmarcado}, no un fracaso. Bajo el patron \textit{Be My Eyes} de Huang et
al.~\cite{huang2025bemyeyes}, el sistema separa \textit{perceiver} y
\textit{reasoner}: los \textit{perceivers} son nuestros modelos densos
(TSViT-pheno, U-TAE) y FarSLIP, que \textit{ven} cada parcela y emiten
observaciones; el \textit{reasoner} es un LLM frontera \textit{frozen}
(Gemini~2.5-pro GA, 1M de contexto, o Qwen~3.5-35B-A3B on-prem para soberania de
datos) que razona sobre esas observaciones en texto plano. El LLM \textbf{no es
un clasificador de pixeles}; el clasificador final es el ensamble heterogeneo.

Por tanto, la descripcion fenologica de FarSLIP no debe juzgarse como cabezal de
clasificacion supervisada (donde su aporte es marginal), sino como \textit{senal
de anclaje} que enriquece la explicacion del \textit{reasoner} sin enviar
tensores ni logits al LLM y sin anadir latencia perceptible. Esta separacion
explica por que el resultado negativo de F2 como clasificador es compatible con
el valor positivo de FarSLIP como componente visión-lenguaje del copiloto.

% ------------------------------------------------------------------------------
\subsection{Caveats explicitos}
\label{sec:farslip_caveats}

\begin{itemize}[leftmargin=*]
    \item \textbf{Reproyeccion cruda.} El lift ortogonal $384\rightarrow768$ es
    un mapeo lineal que preserva norma y angulos, pero \textit{no} es el
    \texttt{CLIPTextModel} nativo: es una aproximacion (caveat
    \texttt{ortho\_proj\_crude\_approx} registrado en los parametros de MLflow).
    La ruta con el codificador de texto nativo queda como trabajo posterior.
    \item \textbf{Tamano de muestra.} La evaluacion FarSLIP-vs-AlphaEarth se
    realiza sobre $n=567$ parcelas (orden $\approx500$--$700$ por fold), con
    validacion cruzada espacial. Las cifras llevan su desviacion estandar.
    \item \textbf{Pureza de parcela.} El filtro \texttt{dominance\_ratio} y la
    agregacion a nivel de parcela mantienen una pureza alta ($\approx98\%$) tras
    descartar los patches sobre-dominados por Meadow; el residual de mezcla se
    asume en la etiqueta de parcela dominante.
    \item \textbf{Saturacion del benchmark.} FarSLIP fiel no supera a AlphaEarth
    como clasificador denso ($-0.0024$); su valor es complementario (ensamble) y
    de transferibilidad, coherente con la transferibilidad espacial limitada
    reportada para embeddings tipo AlphaEarth.
\end{itemize}
